\documentclass[11pt]{article}
\usepackage{mathblatt}
% Präambel des Setzers (werkzeuge/setzer-praeambel.tex, wortgleich); dazu \gteil (Muster T2U)
\makeatletter
\setlength{\mbpfbe}{0mm}% keine Punkte (pruefung.md § 5: nur beim Prüfstein)
\newlength{\gEin}\setlength{\gEin}{0pt}
\let\gSchrift\relax
\renewcommand{\pfkreuzzeile}[1]{\par\noindent\hangindent3.2em\hangafter1\hspace*{1.6em}$\square$\ #1\par}
\newcommand{\gkreuz}[1]{$\square$\ #1\hspace{2em}}
\newcommand{\gkreuzzeile}[2]{\par\noindent\hangindent2.2em\hangafter1\mbox{#1\,$\square$}\ #2\par}% Buchstabe, Option: eine je Zeile (neben Skizze)
% links, was man ansieht, rechts, was man tut; Spaltengrenze überall an derselben Stelle (gemeinsam.md „Satz“)
\newsavebox{\gbox}\newlength{\gLinks}
% v0.10 (Lehrer 13.10.): keine starre Halbierung – die linke Spalte ist so breit wie die Figur (+ 6 mm), die
% Antwort steht direkt daneben; ist die Figur breiter als 0,62 Zeile, steht die Antwort darunter
\newcommand{\gzwei}[2]{\par\smallskip\noindent\sbox{\gbox}{#1}%
  \setlength{\gLinks}{\dimexpr\wd\gbox+6mm\relax}%
  \ifdim\gLinks>0.62\linewidth
    \usebox{\gbox}\par\smallskip\noindent\begin{minipage}{\linewidth}#2\end{minipage}%
  \else
    \begin{minipage}[t]{\gLinks}\vspace{0pt}\usebox{\gbox}\end{minipage}%
    \begin{minipage}[t]{\dimexpr\linewidth-\gLinks\relax}\vspace{0pt}\raggedright #2\end{minipage}%
  \fi\par}
% v0.9 (12.10., pruefung.md § 5 „Marke mit B, Bedingungen rechts“): drittes Argument von pfaufg = Bedingungen
% der Prüfung (v0.10: \pfohnetr = TR durchgestrichen, „FOR“), klein grau im rechten Rand auf der Grundlinie der Nummer, Zeilen mit \\
\newcommand{\gbedrand}[2]{\ifblank{#2}{}{\rlap{\hspace*{\dimexpr#1+2mm\relax}%
  \parbox[t]{15mm}{\raggedright\scriptsize\leavevmode\color{mbgrau}\strut#2}}}}
% Teil eines Bündels (a, b …): Marke links im Rand, Bedingungen rechts im Rand, auf der Zeile des Teils
\newcommand{\gteilrand}[2]{\llap{\mbpfrandmarke{#1}\hspace*{\mbpfnr}}\gbedrand{\linewidth}{#2}}
% Teil eines Bündels: Buchstabe in eigener Spalte, Inhalt eingerückt (\gEin = Einzug, damit die Spaltengrenze von
% \gzwei an derselben Stelle bleibt)
\newenvironment{gbteil}[3]{\par\smallskip\noindent\gteilrand{#2}{#3}\makebox[\mbpfteil][l]{\textbf{#1}}%
  \begin{minipage}[t]{\dimexpr\linewidth-\mbpfteil\relax}\setlength{\gEin}{\mbpfteil}\raggedright\strut\ignorespaces}%
 {\end{minipage}\par}
% Grafik höchstens so breit wie die Spalte (Bündel nebeneinander, Koordinatensystem ohne Karo)
\newsavebox{\gpbox}
\newcommand{\gpassend}[1]{\sbox{\gpbox}{#1}\ifdim\wd\gpbox>\linewidth\typeout{GPASSEND: verkleinert \the\wd\gpbox\space auf \the\linewidth}\resizebox{\linewidth}{!}{\usebox{\gpbox}}\else\usebox{\gpbox}\fi}
% Skizzenfeld (v0.10, Lehrer 13.10.): Karo ab dem linken Rand über die Zeile, „Skizze“ klein grau oben links im Feld
% (wie \pfskizzenfeld der Vorlage); der rechte Rand bleibt den Bedingungen
\newcommand{\gskizze}[1]{\par\smallskip\noindent\begin{tikzpicture}
  \clip (0,0) rectangle ({\linewidth-1mm},{#1*10mm});
  \draw[black!16,line width=0.3pt,step=5mm] (0,0) grid ({\linewidth},{#1*10mm});
  \node[anchor=north west,inner sep=2pt,font=\scriptsize,text=mbgrau,fill=white] at (0,{#1*10mm}) {Skizze};
  \end{tikzpicture}\par}
\RenewDocumentEnvironment{pfaufg}{m m m}{%
  \begin{lrbox}{\mbaufgabebox}\begin{minipage}[t]{\linewidth}%
  \gSchrift\noindent\hspace*{\gEin}\mbpfrandmarke{#1}\gbedrand{\linewidth-\gEin}{#3}%
  \begin{minipage}[t]{\mbpfnr}\strut\textbf{#2}\end{minipage}%
  \begin{minipage}[t]{\dimexpr\linewidth-\gEin-\mbpfnr\relax}\raggedright\strut\ignorespaces}%
 {\end{minipage}%
  \end{minipage}\end{lrbox}%
  \setlength{\mbaufgabehoehe}{\dimexpr\ht\mbaufgabebox+\dp\mbaufgabebox\relax}%
  \ifdefempty{\gkopfnach}{\par\addvspace{5pt}\Needspace*{\mbaufgabehoehe}}{%
    % Stufenkopf, der mit dieser (ersten) Aufgabe zusammenbleibt: Platz für beide verlangen, dann den Kopf setzen
    % (12.10., Quadratische: ein Kopf stand allein unten, die Aufgabe mit Koordinatensystem auf der nächsten Seite)
    \par\Needspace*{\dimexpr\mbaufgabehoehe+40pt\relax}\gkopfnach\gdef\gkopfnach{}\par\addvspace{5pt}}%
  \noindent\usebox{\mbaufgabebox}\par\addvspace{5pt}}
% Stufenkopf vormerken: \gkopfvor{<Kopf>} – die nächste pfaufg setzt ihn, sobald beide zusammen Platz haben
\def\gkopfnach{}
\newcommand{\gkopfvor}[1]{\gdef\gkopfnach{#1}}
% Teilaufgabe mit Skizze links in gzwei: Buchstabe in eigener Spalte, darunter Text und Skizze
\newcommand{\gteil}[2]{\makebox[\mbpfteil][l]{\textbf{#1}}\parbox[t]{\dimexpr0.5\textwidth-\gEin-\mbpfnr-5mm-\mbpfteil\relax}{\raggedright #2}}
% Stufenkopf der Vorlage (mathblatt.sty 2026-10-12: Linie, Abstand, 13 pt, \pfjahre) ohne ihren Vorhalt von
% 14 Zeilen: der Setzer fordert vor dem Kopf 0,16 Seitenhöhe (Kopf + erste Aufgabe), mehr nicht
\setlength{\mbpfstufevorhalt}{0pt}
\makeatother
\newcommand{\karo}[1]{\par\smallskip\noindent\begin{tikzpicture}
  \clip (0,0) rectangle ({\linewidth-1mm},{#1*10mm});
  \draw[black!16,line width=0.3pt,step=5mm] (0,0) grid ({\linewidth},{#1*10mm});
  \end{tikzpicture}\par}
\newcommand{\kk}[1]{$\square$\,#1\hspace{0.9em}}
\newcommand{\linie}[1]{\,{\color{black!45}\rule[-1pt]{#1}{0.4pt}}\,}
\newcommand{\pkt}{\textperiodcentered{}}
% v0.14 (14.10.): Stufenkopf mit fünf Punkten (Option = Jahre für die Bestellkarte) oder „zuletzt 19“
\NewDocumentCommand{\gpunkte}{o m}{#2}
\newcommand{\gzuletzt}[1]{zuletzt #1}
% v0.7 (12.10.): Schülerfrage unter dem Stufenkopf; Ankreuzoption als Bild (Buchstabe und Kästchen darunter)
\newcommand{\gfrage}[1]{\par\nobreak\noindent\hspace*{\mbpfnr}{\small\itshape #1}\par\nobreak\addvspace{2pt}}
\newlength{\gbildbreite}\setlength{\gbildbreite}{0.32\linewidth}% Setzer setzt sie je Reihe (vier Bilder: 0,24)
\newcommand{\gbildoption}[3]{\begin{minipage}[t][#3][b]{\gbildbreite}\centering #2\par\smallskip\mbox{#1\,$\square$}\end{minipage}}
% Lineare (12.10.): Bild mit Buchstabe ohne Kästchen (Diagramm zuordnen); Aussage mit „wahr / falsch“ zum Ankreuzen
\newcommand{\gbildname}[3]{\begin{minipage}[t][#3][b]{\gbildbreite}\centering #2\par\smallskip\mbox{\textbf{#1}}\end{minipage}}
\newcommand{\gwahrfalsch}[2]{\par\noindent\hangindent2.2em\hangafter1\makebox[2.2em][l]{(#1)}#2\hfill\mbox{wahr $\square$\quad falsch $\square$}\par}
\makeatletter\setlength{\mbpfbe}{0mm}\makeatother
\begin{document}

\pfheftstil
\pflaufkopf[P10 \pkt{} mehr]{Netz und Schrägbild}{PG80-5}\pfseitevon
\pfheftkopf{Netz und Schrägbild}{mehr}
\pfstufenzeile{Netz erkennen \pkt{} Schrägbild skizzieren \pkt{} Netz skizzieren}
% 2018-OS-B1i-m1 (Schwere I, rang 4)
\begin{pfaufg}{nach~P10\\’18~\pkt{}~1i}{1.}{}
Kreuze an, welcher Körper zu dem abgebildeten Netz gehört:
\par\smallskip\noindent\gpassend{\begin{tikzpicture}[line width=0.8pt,falt/.style={line width=0.4pt}]\draw (0,0) -- (3.2,0) -- (1.6,2.771) -- cycle;\draw[falt] (1.6,0) -- (0.8,1.386) -- (2.4,1.386) -- cycle;\end{tikzpicture}}\par
\par\smallskip\noindent\pfkreuz{A}{Prisma}\pfkreuz{B}{Pyramide}\pfkreuz{C}{Quader}
\fusshilfe{\mbox{1:~B}}
\end{pfaufg}
% 2017-OS-B1c-m1 (Schwere I, rang 5)
\begin{pfaufg}{nach~P10\\’17~\pkt{}~1c}{2.}{}
Die Abbildung zeigt einen Körper im Schrägbild. Kreuze an, welcher Körper es ist:
\par\smallskip\noindent\gpassend{\begin{tikzpicture}[line width=0.8pt]\draw (0,0)--(1.6,0)--(0.8,1.2)--cycle;\draw (0.8,1.2)--(2.0,1.8)--(2.8,0.6)--(1.6,0);\draw[dashed,line width=0.5pt] (0,0)--(1.2,0.6)--(2.0,1.8) (1.2,0.6)--(2.8,0.6);\end{tikzpicture}}\par
\par\smallskip\noindent\pfkreuz{A}{Pyramide}\pfkreuz{B}{Prisma}\pfkreuz{C}{Quader}
\fusshilfe{\mbox{2:~Prisma}}
\end{pfaufg}
% 2015-OS-K6b-m1 (Schwere I, rang 7)
\begin{pfaufg}{nach~P10\\’15~\pkt{}~6b}{3.}{}
Mia hat ein Modell einer Pyramide mit quadratischer Grundfläche gebaut: Es ist 0,70 m hoch, die Grundkante ist 0,90 m lang. Zeichne in das Schrägbild die Höhe der Pyramide ein. Beschrifte die Skizze mit den gegebenen Maßen.
\par\smallskip\noindent\gpassend{\begin{tikzpicture}[x={(1cm,0cm)},y={(0.5cm,0.5cm)},z={(0cm,1cm)},line width=0.7pt]\draw (0,0,0) -- (3,0,0) -- (3,3,0);\draw (0,0,0) -- (1.5,1.5,2.6) -- (3,0,0);\draw (3,3,0) -- (1.5,1.5,2.6);\draw[dashed,line width=0.5pt] (0,0,0) -- (0,3,0) -- (3,3,0);\draw[dashed,line width=0.5pt] (0,3,0) -- (1.5,1.5,2.6);\end{tikzpicture}}\par
\karo{4}
\fusshilfe{\mbox{3:~0,70 m; 0,90 m}}
\end{pfaufg}
% 2020-OS-K5a-m1 (Schwere II, rang 8)
\begin{pfaufg}{nach~P10\\’20~\pkt{}~5a}{4.}{}
Eine Schachtel hat die Form eines Prismas. Ihre Grundfläche ist ein Rechteck, 11 cm breit und 10 cm tief. Die vordere Kante ist 6 cm hoch, die hintere 14 cm; die Deckfläche verläuft schräg und ist etwa 14 cm breit. Die beiden Seitenflächen sind Trapeze mit zwei rechten Winkeln. Ein Netz ist schon begonnen. Vervollständige das Netz.
\par\smallskip\noindent\gpassend{\begin{tikzpicture}[scale=0.2,line width=0.8pt]\draw[black!14,line width=0.3pt] (0,0) grid (46,39);\draw (1,14) rectangle (14.6,24);\draw (14.6,14) rectangle (20.6,24);\draw (20.6,14) rectangle (31.6,24);\draw (20.6,24)--(20.6,30)--(31.6,38)--(31.6,24);\node[font=\scriptsize,mbgrau,above right,inner sep=1pt] at (0,0) {1 Kästchen $\mathrel{\widehat{=}}$ 1 cm};\end{tikzpicture}}\par
\karo{3}
\fusshilfe{\mbox{4:~Rückfläche 14 cm × 10 cm rechts an die Grundfläche; zweites Trapez unten an die Grundfläche}}
\end{pfaufg}
% 2014-OS-K5b-m1 (Schwere II, rang 10)
\begin{pfaufg}{nach~P10\\’14~\pkt{}~5b}{5.}{}
Für ein Café wird eine Kühlvitrine in Form eines Prismas gebaut. Die Grundfläche ist ein gleichschenkliges Trapez mit den parallelen Seiten 108 cm und 84 cm, den Schenkeln 37 cm und dem Flächeninhalt 3 360 cm². Die Vitrine ist 150 cm hoch. Ein Netz der Vitrine ist schon begonnen. Vervollständige das Netz. Gib die fehlenden Maße a, b und d an.
\par\smallskip\noindent\gpassend{\begin{tikzpicture}[line width=0.8pt,font=\small]\draw (0,0) rectangle (0.74,3.0);\draw (0.74,0) rectangle (2.42,3.0);\draw (0.74,3.0)--(0.5,3.7)--(2.66,3.7)--(2.42,3.0);\draw (0.74,0)--(0.5,-0.7)--(2.66,-0.7)--(2.42,0);\node[above] at (0.37,3.0) {$b$};\node[below] at (1.58,3.0) {$c$};\node[left] at (0,1.5) {$d$};\node[below] at (1.58,-0.7) {$a$};\end{tikzpicture}}\par
\karo{3}
\fusshilfe{\mbox{5:~Netz um Rechtecke 37 cm × 150 cm und 108 cm × 150 cm ergänzt; 150 cm}}
\end{pfaufg}
\par\vfill\noindent{\scriptsize\color{mbgrau}Originale: P10 2018 \pkt{} 1i, P10 2017 \pkt{} 1c, P10 2015 \pkt{} 6b, P10 2020 \pkt{} 5a, P10 2014 \pkt{} 5b\par}
\end{document}